% Complete Spanish content, retained from the previous edition.
\chapter{Qué necesitas}


\begin{table}[H]\centering\small
\begin{tabularx}{\linewidth}{@{}lL@{}}
\toprule
\textbf{Elemento} & \textbf{Requisito}\\
\midrule
Ordenador & MSX2, MSX2+ o MSX turbo R\\
Memoria & 64 KB de RAM y 128 KB de VRAM\\
Vídeo & Vídeo propio del MSX2; SCREEN~5, 50 o 60~Hz.
No necesita una tarjeta de vídeo adicional\\
Sonido & PSG interno; opcionalmente OPL4 (MoonSound o compatible)\\
Juego & Cartucho MegaROM de 2 MiB, tipo ASCII16\\
Conexión & Una ranura libre para el juego; otra si se utiliza MoonSound\\
Pantalla & Monitor o televisor conectado a la salida de vídeo del MSX\\
Control & Teclado; opcionalmente joystick MSX en puerto 1 o ratón MSX
en puerto 1 o 2\\
\bottomrule
\end{tabularx}
\caption{Equipo necesario.}
\end{table}

No necesitas cinta ni disco. Para conservar a los héroes entre sesiones,
apunta sus códigos de guardado (sección~\ref{sec:codigos}).

\chapter{Puesta en marcha}


\begin{enumerate}
\item Comprueba que el ordenador está \textbf{apagado}.
\item Inserta el cartucho del juego en una ranura libre y presiónalo
      hasta que quede firme. Si utilizas MoonSound, insértalo en otra
      ranura libre, también con el ordenador apagado.
\item Si vas a usar ratón o joystick, conéctalo ahora (véase el capítulo
      \ref{cap:controles}).
\item Conecta la pantalla a la salida de vídeo del MSX.
\item Enciende el ordenador. El juego arranca solo.
\end{enumerate}

\aviso{¡Atención! Cuidado con el cartucho}{Inserta y extrae el cartucho
\textbf{siempre con el ordenador apagado}. No lo introduzcas ni lo retires
nunca con el ordenador encendido: podrías dañar tanto el cartucho como el
propio ordenador, y además perderías la partida en curso.}

\chapter{Del encendido al juego}


Al encender el ordenador se suceden, por este orden:

\begin{enumerate}
\item \textbf{El logotipo de DiHalt y su fanfarria musical}. Esta secuencia dura
      unos 9,4 segundos con el PSG interno, incluida la pausa final, o
      10 segundos con MoonSound. Avanza automáticamente y no se puede saltar.
\item \textbf{La ilustración de apertura}. Pulsa \textbf{ESPACIO} o el botón del
      joystick, o haz clic con el botón izquierdo del ratón, para continuar.
\item \textbf{La selección de idioma}. Elige ENGLISH o ESPAÑOL con arriba/abajo y confirma con \textbf{ESPACIO}, \textbf{RETURN} o disparo. Español aparece seleccionado por defecto. La intro y el juego utilizarán ese idioma hasta reiniciar. Los códigos de héroe sirven en ambos idiomas.
\item \textbf{La introducción animada}, en la que Mentor cuenta la historia
      de Morcar. Puedes saltarla con
      \textbf{ESPACIO}, \textbf{ESC} o con el botón derecho del ratón.
\item \textbf{El menú principal}.
\end{enumerate}

\begin{figure}[H]\centering
\includegraphics[width=0.8\linewidth,height=0.6\linewidth]{opening16}
\caption{La ilustración de apertura tal como la muestra el MSX2, con sus
dieciséis colores.}
\end{figure}

\pantalla[0.62]{language_selector}{Selecciona el idioma antes de comenzar la introducción.}
\subsection{Idioma}
El cartucho contiene ambos idiomas. La elección se aplica a toda la introducción, los menús, los mensajes y las veinticuatro aventuras. Para elegir de nuevo, reinicia el juego; la selección de idioma no se guarda. Anota los códigos de tus héroes antes de reiniciar si quieres conservarlos. Los códigos de héroe son compatibles entre ambos idiomas y las reglas no cambian.

\subsection{Lo que no se conserva al apagar}


El cartucho no graba datos: al apagar el ordenador se pierde la partida.
Para no perder a tus héroes, el juego te da de cada uno un \textbf{código de
16 letras} que guarda su nombre, su oro, sus puntos de cuerpo, sus armas, su
armadura y los tesoros de aventura que ha encontrado. Apúntalo antes de
apagar y tecléalo otro día para recuperar al héroe tal como estaba (se
explica en la página~\pageref{sec:codigos}).

% ================================================================
\chapter{Controles}
\label{cap:controles}

Todo el juego se maneja con un \textbf{puntero} en forma de flecha y un
único \textbf{botón de acción}. Puedes moverlo con cualquiera de estos
dispositivos, sin necesidad de configurar nada:

\begin{table}[H]\centering\small
\begin{tabularx}{\linewidth}{@{}lLL@{}}
\toprule
\textbf{Dispositivo} & \textbf{Mover el puntero} & \textbf{Acción}\\
\midrule
Ratón MSX (puerto 1 o 2) & Desplazar el ratón & Botón izquierdo\\
Teclado & Teclas del cursor & ESPACIO\\
Joystick (puerto 1) & Palanca & Botón de disparo\\
\bottomrule
\end{tabularx}
\caption{Dispositivos de control.}
\end{table}

\subsection{Ratón}

El juego detecta por sí solo si hay un ratón MSX conectado, en el puerto 1
o en el 2, y lo utiliza junto con el teclado. El botón izquierdo equivale a
la acción; el botón derecho sólo sirve para saltar la introducción. Con el
ratón el juego resulta mucho más ágil: basta señalar un icono, una casilla
o una opción y pulsar.

\subsection{Teclado y joystick}

Sin ratón se juega exactamente igual: las teclas del cursor (o la palanca
del joystick conectado al puerto 1) desplazan el puntero por la pantalla, y
ESPACIO (o el disparo) pulsa sobre lo que señala.

\subsection{En los menús}

En los menús hay siempre una línea resaltada. Con el teclado o el joystick,
las teclas \textbf{arriba} y \textbf{abajo} cambian la línea resaltada y
ESPACIO la elige. Sin ratón, la flecha se oculta en los menús; durante la partida vuelve a
aparecer para manejar los iconos. Con el ratón basta llevar el puntero sobre la opción
deseada ---el resalte lo sigue--- y pulsar el botón izquierdo; una
pulsación fuera de las opciones no hace nada.

\begin{table}[H]\centering\small
\begin{tabularx}{\linewidth}{@{}lL@{}}
\toprule
\textbf{Tecla} & \textbf{Función}\\
\midrule
Cursores & Mover el puntero o la línea resaltada\\
ESPACIO & Acción: elegir, pulsar un icono, detener el número de movimiento
          o saltar la introducción\\
ESC & Saltar la introducción\\
\bottomrule
\end{tabularx}
\caption{Resumen de teclas.}
\end{table}

% ================================================================
\chapter{Introducción}


\textit{HeroQuest} es un juego de aventuras y fantasía por turnos, basado en
el célebre juego de tablero. Hasta cuatro héroes ---un bárbaro, un enano, un
duende y un hechicero--- se adentran en mazmorras repletas de monstruos,
trampas y tesoros para cumplir las misiones que les encomienda Mentor, el
guardián del libro Loretome. Enfrente tienen a Morcar, el Señor del Caos, a
quien da vida el ordenador: él mueve a los monstruos y decide sus ataques,
de modo que todos los jugadores pueden colaborar en el mismo bando.

Esta versión para MSX2 reúne en un solo cartucho las catorce aventuras del
juego original de Gremlin Graphics y las diez de su expansión, \textit{El
retorno del Señor Brujo}: veinticuatro aventuras en total, con una introducción animada, control por ratón, teclado o
joystick y música digital si el ordenador dispone de un cartucho OPL4.

Este manual corresponde a la versión \versionjuego{}, cuyo número aparece
en el menú principal, debajo de los créditos de DiHalt.

Pueden jugar de una a cuatro personas. Con cuatro jugadores cada uno lleva
un héroe; con menos, cada jugador puede dirigir varios, y un jugador
solitario puede aventurarse con un único personaje o con el grupo completo.

Conviene leer este manual antes de empezar: las reglas son sencillas, pero
la mazmorra no perdona a quien entra en ella sin conocerlas.

\pantalla[0.8]{sala_monstruos}{El bárbaro entra en una sala ocupada por los
siervos de Morcar.}

\subsection{Las crónicas de Loretome}

<<Escuchad con atención las palabras de Mentor, guardián de Loretome, pues
voy a hablaros de tiempos pasados, de días oscuros en los que el Imperio se
salvó contra toda esperanza. Temo que la oscuridad esté a punto de
regresar.>>

Las legiones de Morcar, Señor del Caos, lo habían arrasado todo a su paso.
Ante el Estandarte Negro y las hordas del Caos, hasta los guerreros más
valientes del Emperador dieron media vuelta y huyeron. La tierra quedó
devastada y los hombres perdieron la esperanza.

Entonces llegó de las Tierras Fronterizas un poderoso príncipe guerrero:
Rogar el Bárbaro. Lucía en la frente un diamante resplandeciente, la
Estrella del Oeste, la misma joya que portaron los antiguos reyes de la
leyenda. Con él volvió la esperanza, y los hombres abandonaron sus
escondites en colinas y bosques para unirse a su estandarte. Otros grandes
héroes se le sumaron: Durgin, el intrépido guerrero enano de las Montañas
del Fin del Mundo; Ladril, el guerrero mago de los elfos, llegado de la
lejana Athelorn; y Telor, el hechicero cuya magia salvaría a Rogar en más de
una ocasión.

Durante años Rogar adiestró a su ejército. Evitó la batalla campal contra
el general de Morcar hasta que todo estuvo dispuesto, pero no dejó de
hostigar sus líneas de suministro ni de acabar con orcos y goblins. Y llegó
el día esperado. Acampado en los pasos altos, Ladril divisó a lo lejos las
huestes negras y pidió a Durgin que tocara la llamada a las armas con su
poderoso cuerno. Los ejércitos de Rogar cayeron sobre el enemigo por dos
flancos. Muchas criaturas inmundas y muchos hombres buenos perecieron
aquella jornada; pero, al apagarse la luz del día, fue la Oscuridad la que
abandonó el campo.

La victoria, sin embargo, no fue completa. Morcar y su general escaparon
más allá del Mar de las Garras, y todavía hoy traman su venganza. Pronto
sus planes estarán maduros y el Imperio necesitará un nuevo Rogar. ¿Dónde
están los héroes capaces de igualarlo? Mucho tenéis que aprender para
llegar a ser tan grandes como él y sus compañeros.

Mentor os ayudará en lo que pueda. El libro que lleva consigo, Loretome,
fue escrito cuando comenzó el tiempo, y en sus páginas incontables está
registrado todo lo que fue y todo lo que será. A través de él puede
guiaros; pero no le está permitido intervenir, pues entonces un mal mayor
caería sobre el mundo y el Caos triunfaría para siempre.

% ================================================================
\chapter{Sonido}



El juego reconoce automáticamente el equipo de sonido disponible; no hay
nada que configurar. Al arrancar, el logotipo de DiHalt tiene su propia
fanfarria musical, que suena por MoonSound si está presente o por el PSG
interno en caso contrario. Después del logotipo, la introducción, los menús
y las aventuras utilizan la música que se describe a continuación.

\subsection{Con el sonido interno}

Con el sonido interno del MSX suenan la música
del juego original ---una melodía para los menús y la introducción y otra
para las aventuras--- y sus efectos: los pasos, las puertas, los golpes y
las heridas. La música y los efectos no suenan a la vez: mientras la música está activada no se oyen los efectos, y al
desactivarla (véase más abajo) el juego pasa a reproducirlos.

\subsection{Con un cartucho OPL4 (MoonSound)}

Si al arrancar el juego encuentra un OPL4, la música pasa a sonar por él:
es música digital, tomada de los módulos de la versión de Amiga, con una
melodía para la introducción y los menús y otra distinta durante las
aventuras. También se escuchan los efectos de sonido del juego, de modo que música
y efectos suenan a la vez.

\subsection{Activar y desactivar la música}

En la pantalla de opciones (opción \pant{Opciones} del menú principal),
la opción \pant{Música} la activa o la desactiva. Sin OPL4, desactivarla
es la manera de jugar con efectos de sonido.

% ================================================================
\chapter{La introducción}
\label{cap:intro}

Tras la ilustración de apertura y la selección de idioma comienza la introducción animada, con los textos en el idioma elegido. Cae la
tarde sobre la ciudad y, en su estudio, a la luz de un fuego que calienta
pero apenas alumbra, Mentor reúne a sus discípulos. Entre sombras que dejan
adivinar estanterías repletas de libros y pergaminos, el viejo mago se
acerca a la chimenea y les habla.

\dospantallas{intro_titulo}{intro_ciudad}{El título y la primera escena de
la introducción.}

Su entrenamiento ha terminado, les dice, pero todavía no son héroes: antes
han de ponerse a prueba. Y para que sepan a qué se enfrentan les cuenta la
historia de Morcar.

Hace muchos siglos Morcar fue su aprendiz. Trabajaba duro y aprendía
deprisa, pero la impaciencia lo devoraba: quería conocer magias cada vez
más poderosas. Mentor le advirtió de los peligros y le pidió paciencia,
pues con el tiempo llegaría a ser un gran mago. Morcar no quiso esperar.
Noche tras noche entró a escondidas en el estudio de su maestro para leer
sus libros de hechizos y, cuando hubo aprendido los grandes secretos que
guardaban, huyó.

\dospantallas{intro_mentor}{intro_texto}{Mentor relata la traición de su
aprendiz; los textos de la introducción están en castellano.}

Cuando Mentor dio con él lo encontró muy cambiado: había jurado lealtad a
los grandes poderes del Caos. Intentó hacerle entrar en razón, pero Morcar
se rió en su cara y desató contra él un hechizo terrible. Lucharon durante
días. Morcar contaba con aliados más fuertes que el propio Mentor, que no
pudo derrotarlo; y cuando ambos estaban ya exhaustos, el traidor escapó y
buscó refugio en los yermos del Caos, al norte.

Allí curó sus heridas, perfeccionó sus artes y conjuró poderes antiguos con
los que derrocar al Imperio; poderes que le servirán mientras convenga a
sus oscuros propósitos, y que acabarán por destruirlo. Desde entonces
Mentor vigila a Morcar y mide la fuerza de su magia, porque lo que ha
convocado acabaría con todos si esa vigilancia se relajara. Sus legiones ya
amenazaron al Imperio una vez, y fue Rogar quien las derrotó. Ahora están
de nuevo en marcha y han atacado las fronteras.

\dospantallas{intro_estudio}{intro_sala}{Otras dos escenas de la
introducción.}

El Imperio vuelve a necesitar héroes, y para eso han sido entrenados. Cada
uno deberá superar tres pruebas; quien lo consiga será aclamado campeón del
reino y nombrado caballero imperial, y estará de verdad en el camino de los
héroes. <<Hablaré con vosotros otra vez cuando volváis\ldots{} si volvéis.>>

La introducción se muestra una sola vez, al encender el ordenador. Puedes
saltarla en cualquier momento con \textbf{ESPACIO}, \textbf{ESC} o con el
botón derecho del ratón.

% ================================================================
\chapter{Los menús}
\label{cap:menus}

\subsection{El menú principal}

\pantalla[0.62]{menu_principal}{El menú principal de la versión \versionjuego.}

El menú principal ofrece cuatro opciones:

\begin{description}
\item[\pant{Jugar}] Comienza una aventura. Debería ser lo último que elijas,
  una vez preparados los personajes.
\item[\pant{Comprar equipo}] Lleva a la tienda, donde los héroes gastan su
  oro en armas y armaduras.
\item[\pant{Crear personaje}] Permite decidir qué héroes participan y darles
  nombre.
\item[\pant{Opciones}] Ajustes de música y tipo de letra.
\end{description}

\subsection{Crear personaje}

\pantalla[0.62]{crear_personaje}{La pantalla de personajes: a la izquierda
la clase de cada héroe y a la derecha su nombre.}

La pantalla \pant{Elige jugador} muestra a los cuatro héroes. A la derecha
de cada uno aparece su nombre ---de fábrica, Sigmar el bárbaro, Grugni el
enano, Ladril el duende y Zoltan el hechicero--- o bien la palabra
\pant{Muerto}. Un personaje muerto no toma parte en la siguiente aventura.

\begin{itemize}
\item Con \textbf{arriba} y \textbf{abajo} (o señalando con el ratón) se
  resalta un héroe.
\item Con \textbf{derecha} se alterna entre el nombre y \pant{Muerto}. Así
  se retira del grupo a un héroe que no vaya a jugar, y así se crea un
  personaje nuevo en el lugar de uno caído: el nuevo empieza con sus puntos
  al completo, pero sin oro ni equipo.
\item Pulsando la acción sobre un héroe activo se abre su menú, con las
  opciones \pant{Nombre jugador}, \pant{Cargar} y \pant{Guardar}.
\item \pant{Volver} devuelve al menú principal.
\end{itemize}

\textbf{Nombre jugador} muestra un abecedario en la parte inferior de la
pantalla: elige las letras una a una, hasta un máximo de seis, y termina con
la opción de fin.

\pant{Cargar} y \pant{Guardar} servían en el juego original para llevar los
personajes a una cinta. En el cartucho trabajan con \textbf{códigos}, como
se explica a continuación.

Si todos los héroes han muerto, el juego lo advierte y pide crear
personajes nuevos antes de volver a jugar.

\subsection{Los códigos de héroe}\label{sec:codigos}

El cartucho no puede grabar, así que cada héroe se conserva con un código
de 16 letras: 14 de datos y dos de comprobación. Los códigos antiguos de
14 letras ya no se aceptan. El código contiene:

\begin{itemize}
\item el nombre del héroe (sus seis letras);
\item su oro;
\item sus puntos de cuerpo en ese momento;
\item las armas y la armadura que ha comprado, y la caja de herramientas;
\item los tesoros de aventura que ha encontrado.
\end{itemize}

No guarda las pociones ni los hechizos, que se pierden al terminar cada
aventura, ni qué arma llevaba en la mano: al recuperar al héroe hay que
volver a elegirla en el inventario.

\dospantallas{codigo_guardar}{codigo_cargar}{A la izquierda, \pant{Guardar}
muestra el código del héroe. A la derecha, \pant{Cargar} lo pide letra a
letra.}

\subsubsection*{Guardar: ver el código}

En \pant{Crear personaje}, elige al héroe, pulsa la acción y escoge
\pant{Guardar}. La pantalla muestra su código en tres grupos de letras,
por ejemplo \texttt{QGYNT ZDWQL PYZGTI}. Apúntalo tal cual y pulsa la acción
para volver. Conviene hacerlo con los cuatro héroes al terminar una aventura
y después de pasar por la tienda: el código refleja al héroe en el momento
en que se pide.

\subsubsection*{Cargar: teclear el código}

Elige al héroe del mismo tipo (un código de bárbaro solo sirve para el
bárbaro), pulsa la acción y escoge \pant{Cargar}. Las letras se pueden
introducir de tres maneras:

\begin{itemize}
\item tecleándolas directamente en el teclado;
\item señalando cada letra del abecedario con el ratón y pulsando el botón;
\item moviendo el cursor del abecedario con \textbf{izquierda} y
  \textbf{derecha} y pulsando la acción.
\end{itemize}

La opción \pant{Fin} borra la última letra; con la línea vacía, devuelve a
la pantalla de personajes sin cambiar nada. Al completar la letra número 16
el héroe aparece con su nombre, su oro y su equipo. Si alguna letra está
mal, el juego avisa con \pant{Código incorrecto} y deja volver a
escribirlo: es casi imposible que un código mal copiado sea aceptado.

\aviso{Atención}{Cargar un código sustituye al héroe que hubiera en ese
puesto. Si quieres conservarlo, apunta antes su código.}

\subsection{Comprar equipo}

\pantalla[0.7]{tienda}{La tienda. Arriba, los compradores; abajo, las armas
y armaduras con su valor.}

A lo largo de las aventuras los héroes encuentran oro y cobran recompensas.
En la tienda pueden gastarlo en equipo mejor, cuyas ventajas se aplican
automáticamente.

\begin{enumerate}
\item Elige en la lista superior al héroe que va a comprar. Su oro aparece
  en pantalla.
\item Pulsa sobre el artículo deseado. Si el héroe tiene oro suficiente, y
  su clase puede usarlo, el artículo pasa a ser suyo.
\item \pant{Volver}, en la lista de artículos, cierra la compra de ese héroe;
  \pant{Volver}, en la lista de héroes, devuelve al menú principal.
\end{enumerate}

\begin{table}[H]\centering\small
\begin{tabularx}{\linewidth}{@{}lrL@{}}
\toprule
\textbf{Arma} & \textbf{Valor} & \textbf{Efecto}\\
\midrule
\pant{Bastón}     & 100 & 2 dados de ataque. Ataca también en diagonal.\\
\pant{Daga}       & 150 & 2 dados de ataque. Ataca también en diagonal.\\
\pant{Espada}     & 250 & 3 dados de ataque.\\
\pant{Gran hacha} & 400 & 4 dados de ataque. Impide usar el escudo.\\
\pant{Ballesta}   & 350 & 3 dados de ataque a distancia. No puede usarse
                          contra un enemigo adyacente.\\
\pant{Hacha}      & 150 & 2 dados de ataque. Se puede arrojar, pero
                          entonces se pierde.\\
\pant{Lanza}      & 150 & 2 dados de ataque. Ataca en diagonal. Se puede
                          arrojar, pero entonces se pierde.\\
\bottomrule
\end{tabularx}
\caption{Armas a la venta.}
\end{table}

\begin{table}[H]\centering\small
\begin{tabularx}{\linewidth}{@{}lrL@{}}
\toprule
\textbf{Armadura} & \textbf{Valor} & \textbf{Efecto}\\
\midrule
\pant{Escudo}   & 100 & Un dado más de defensa.\\
\pant{Casco}    & 120 & Un dado más de defensa.\\
\pant{Malla}    & 450 & Cota de malla: 3 dados de defensa.\\
\pant{Armadura} & 850 & Armadura de placas: 4 dados de defensa, pero su
                        peso reduce el movimiento a un solo dado.\\
\pant{Caja}     & 250 & Caja de herramientas: permite desactivar las
                        trampas descubiertas. El enano ya tiene una.\\
\bottomrule
\end{tabularx}
\caption{Armaduras y equipo a la venta.}
\end{table}

El \textbf{hechicero} sólo puede comprar el bastón y la caja de
herramientas: ni las demás armas ni las armaduras son para él.

\subsection{Opciones}

\pantalla[0.62]{controles}{La pantalla de opciones.}

Se maneja como el menú principal: \textbf{arriba} y \textbf{abajo} mueven
la barra y la acción elige, o se señala la opción con el ratón y se pulsa
el botón.

\begin{description}
\item[\pant{Música}] Activa o desactiva la música (\pant{Sí} o
  \pant{No}).
\item[\pant{Letra}] Cambia el tipo de letra de todo el juego entre la
  gótica, que es la de fábrica, y otra normal, más fácil de leer.
\item[\pant{Volver}] Vuelve al menú principal.
\end{description}

\subsection{Jugar: hechizos y elección de aventura}

Al elegir \pant{Jugar}, y antes de escoger la aventura, los personajes que
usan magia se reparten los hechizos. Hay cuatro grupos ---\pant{Aire},
\pant{Fuego}, \pant{Tierra} y \pant{Agua}---, cada uno con tres hechizos.

\pantalla[0.7]{hechizos}{Elección de los grupos de hechizos.}

El aviso superior indica quién debe elegir: \pant{Hechicero selecciona
hechizos} o \pant{Duende selecciona hechizos}.

El hechicero elige primero un grupo; a continuación el duende elige el
suyo, y los dos grupos restantes quedan para el hechicero. Si sólo juega
uno de los dos, elige él solo: el duende un grupo, el hechicero tres.

Después aparece la lista de aventuras, que tiene dos páginas:

\dospantallas{aventuras1}{aventuras2}{Las dos páginas de la lista de
aventuras: el juego original y \textit{El retorno del Señor Brujo}.}

\begin{itemize}
\item La primera contiene las \textbf{catorce aventuras del juego
  original}. Su última línea, \pant{Más}, pasa a la segunda página.
\item La segunda contiene las \textbf{diez aventuras de la expansión}. Su
  última línea, \pant{Volver}, regresa a la primera.
\end{itemize}

Con el ratón, la misión que señalas queda resaltada directamente.

Al elegir una aventura aparecen las palabras con que Mentor la presenta.
Léelas con atención, porque explican qué hay que conseguir, y pulsa la
acción para entrar en la mazmorra. Lo mejor es jugar las aventuras en
orden: cada una es algo más difícil que la anterior y la historia avanza
de una a otra.

\pantalla[0.7]{pergamino1}{Mentor presenta la primera aventura.}

% ================================================================
\chapter{Cómo se juega}


\subsection{La pantalla de juego}

La mazmorra se ve sala a sala, en perspectiva isométrica, dentro de un
marco con los datos del héroe que tiene el turno y los iconos de las
acciones.

Arriba a la izquierda, el corazón rojo indica los puntos de cuerpo; la
bolsa dorada, el oro, y la huella cian, los movimientos disponibles. Los
avisos y las opciones de búsqueda se muestran sobre pergaminos con bordes
y texturas que varían según el color de fondo de la sala.

\pantalla[0.80]{pantalla_juego}{La pantalla de juego.}

\begin{table}[H]\centering\small
\begin{tabularx}{\linewidth}{@{}rL@{}}
\toprule
1 & Nombre del héroe que tiene el turno.\\
2 & Corazón rojo: puntos de cuerpo. Si llegan a cero, el héroe muere.\\
3 & Puntos de mente.\\
4 & Bolsa dorada: oro que lleva el héroe.\\
5 & Huella cian: movimientos que le quedan en este turno.\\
6 & Vista de la sala o del pasillo. Pulsando sobre una casilla, el héroe
    camina hasta ella.\\
7 & Siguiente jugador: termina el turno.\\
8 & Llave: abre la puerta que el héroe tiene delante.\\
9 & Flechas de dirección: mueven al héroe una casilla.\\
10 & Ojo: buscar (trampas y puertas secretas, o tesoros).\\
11 & Bolsa: inventario del héroe (pociones, armas, armadura, hechizos y
     tesoros de aventura).\\
12 & Espada: atacar.\\
13 & Pergamino: mapa de la mazmorra.\\
\bottomrule
\end{tabularx}
\caption{Elementos de la pantalla de juego.}
\end{table}

\subsection{El turno}

Los héroes juegan por orden, uno tras otro; el nombre del que tiene el
turno aparece arriba a la izquierda. Cuando todos han jugado le toca a
Morcar, y después se vuelve a empezar.

En su turno, cada héroe puede hacer dos cosas: \textbf{moverse} y realizar
\textbf{una acción} (atacar, lanzar un hechizo o buscar). Puede moverse
primero y actuar después, o actuar primero y moverse después; lo que no
puede es moverse, actuar y seguir moviéndose. Nadie está obligado a actuar
ni a agotar su movimiento. Cuando hayas terminado, pulsa el icono de
siguiente jugador.

\subsection{El movimiento}

Al comenzar el turno, el número de movimientos gira sin parar: hace las
veces de los dados del juego de tablero. \textbf{Pulsa la acción para
detenerlo}; la cifra que quede es el número de casillas que el héroe puede
recorrer en ese turno. El juego ya tiene en cuenta todo lo que afecta a la
tirada: los héroes tiran normalmente dos dados, y sólo uno si visten
armadura de placas o cargan con algo muy pesado.

Hay dos maneras de mover al héroe:

\begin{itemize}
\item con las \textbf{flechas de dirección}, casilla a casilla;
\item pulsando en la vista sobre la \textbf{casilla de destino}: si está en
  línea recta y quedan movimientos suficientes, el héroe camina hasta ella.
\end{itemize}

Reglas del movimiento:

\begin{itemize}
\item No se puede mover en diagonal.
\item No se puede terminar en una casilla ocupada por un monstruo o un
  mueble, ni atravesar muros o escombros (\pant{Bloqueada}).
\item Dos héroes pueden coincidir en una casilla, pero ninguno de los dos
  puede atacar desde ella.
\item Los monstruos tienen un movimiento fijo, sin tirada.
\end{itemize}

\subsubsection*{Puertas}

Héroes y monstruos sólo pueden entrar y salir de las salas por las puertas
abiertas. Los monstruos no saben abrirlas; los héroes sí, de dos formas:
colocándose ante la puerta y pulsando el icono de la \textbf{llave}, o
simplemente caminando a través de ella, pues el juego la abre solo. Abrir
una puerta no gasta movimiento, y la puerta queda abierta para el resto de
la aventura. En un mismo turno no se puede volver a cruzar la puerta por la
que se acaba de pasar.

\subsubsection*{La salida}

En casi todas las aventuras los héroes empiezan junto a una escalera de
caracol, que es también la \pant{Salida}. Un héroe que pisa la escalera
abandona la mazmorra. La aventura termina cuando ningún héroe queda dentro.

\pantalla[0.78]{sala_inicio}{Los cuatro héroes junto a la escalera, al
comienzo de una aventura.}

\subsection{Buscar}

En lugar de atacar o lanzar un hechizo, el héroe puede buscar, antes o
después de moverse. Al pulsar el ojo, el juego pregunta qué se busca:

\pantalla[0.78]{buscar}{El menú de búsqueda: trampas y puertas secretas, o
tesoros.}

\begin{description}
\item[\pant{Revisas trampas, puertas}] Examina la sala o el tramo de pasillo
  en que está el héroe. Las puertas secretas halladas quedan visibles y
  abiertas para el resto de la aventura, y las trampas quedan al
  descubierto.
\item[\pant{Buscas tesoros}] Registra la sala. Puede aparecer oro, una
  poción, el tesoro especial de la aventura\ldots{} o algo peor.
\end{description}

No se puede buscar tesoros mientras haya monstruos en la misma sala.
Los monstruos nunca buscan.

\subsubsection*{Tesoros}

En algunas salas se esconde un tesoro concreto, del que suele hablar la
presentación de la aventura. En las demás, la búsqueda es cuestión de
suerte: lo más habitual es encontrar unas monedas (de 10 a 100), y a veces
una poción; pero también puede no haber nada, saltar una trampa que cuesta
un punto de cuerpo, o aparecer un \textbf{monstruo errante} que ha seguido
al héroe y lo ataca en el acto. Cada sala sólo puede registrarse una vez.

Algunos cofres están protegidos por una trampa. Si antes de buscar tesoros
se revisa la sala en busca de trampas, la del cofre queda desactivada.

\subsubsection*{Trampas}

Las trampas no se ven hasta que un héroe las descubre buscando o cae en
ellas. Los monstruos no las activan.

\begin{description}
\item[Foso] Quien cae en él pierde un punto de cuerpo. El foso queda a la
  vista como obstáculo (\pant{Trampa}); se puede intentar saltar, con el
  riesgo de tropezar, caer y perder otro punto.
\item[Desprendimiento] Al pasar, el techo se viene abajo: el héroe puede
  resultar herido y una roca bloquea el paso para siempre. Cuidado: puede
  cortar el camino de vuelta.
\item[Lanza] Una lanza sale disparada del muro contra el primer héroe que
  pisa la casilla. Puede herirlo o fallar; en cualquier caso sólo se
  dispara una vez. Si se descubre al buscar, queda inutilizada.
\end{description}

Quien activa una trampa no puede hacer nada más hasta su turno siguiente.

El \textbf{enano} lleva siempre su caja de herramientas, y con ella puede
desactivar una trampa descubierta que tenga junto a él: abre el inventario
y pulsa sobre la caja. Los demás héroes pueden hacer lo mismo si compran
una caja en la tienda. Desactivar una trampa tiene su riesgo: si sale mal,
la trampa salta y cuesta un punto de cuerpo.

\subsection{El combate}

Para atacar a un monstruo hay que estar \textbf{junto a él}: delante,
detrás o a un lado. No se puede atacar en diagonal, salvo con las armas que
lo permiten (bastón, daga y lanza). Con la ballesta, o arrojando el hacha o
la lanza, se puede atacar a distancia a cualquier enemigo que esté a la
vista y no sea adyacente.

Para atacar:

\begin{enumerate}
\item Pulsa el icono de la \textbf{espada}. La pantalla cambia al mapa.
\item Pulsa sobre el enemigo elegido. (En la parte inferior, \pant{Arma}
  abre el inventario para cambiar de arma y \pant{Cancelar} vuelve a la
  sala sin atacar.)
\item Aparece la pantalla de combate, con el resultado.
\end{enumerate}

\dospantallas{mapa_ataque}{combate_victoria}{Elección del enemigo sobre el
mapa y resultado de un ataque afortunado.}

El ordenador tira los dados por los dos bandos. Son dados especiales de
combate, de seis caras: tres tienen una \textbf{calavera}, dos un
\textbf{escudo de héroe} y una un \textbf{escudo de monstruo}.

\begin{itemize}
\item El atacante tira tantos dados como indique su ataque. Cada calavera
  es un impacto.
\item El defensor tira sus dados de defensa. Cada escudo \emph{de su bando}
  anula un impacto: los héroes se defienden con dos caras de cada dado; los
  monstruos, sólo con una.
\item Cada impacto no anulado resta un punto de cuerpo.
\end{itemize}

La gran mayoría de los monstruos tiene un solo punto de cuerpo, así que
caen con el primer impacto que no detengan. Un héroe muere cuando sus
puntos de cuerpo llegan a cero.

\begin{table}[H]\centering\small
\begin{tabular}{@{}lcc@{}}
\toprule
\textbf{Héroe} & \textbf{Dados de ataque} & \textbf{Dados de defensa}\\
\midrule
Bárbaro   & 3 & 2\\
Enano     & 2 & 2\\
Duende    & 2 & 2\\
Hechicero & 1 & 2\\
\bottomrule
\end{tabular}
\caption{Dados de combate de los héroes sin equipo adicional.}
\end{table}

Las armas y armaduras cambian estas cifras, como se indica en las tablas de
la tienda. Si un héroe ha derrotado a su adversario antes de moverse,
todavía puede hacerlo. Asegúrate de estar bien colocado antes de pulsar la
espada: si el enemigo elegido queda fuera de alcance, el juego avisa de que
estás demasiado lejos y la acción del turno se pierde.

\dospantallas{combate_bloqueado}{combate_herido}{El turno de Morcar: un
ataque enemigo detenido y otro que hiere al bárbaro.}

\begin{figure}[H]\centering
\includegraphics[width=0.32\linewidth]{lamina_0}\hfill
\includegraphics[width=0.32\linewidth]{lamina_1}\hfill
\includegraphics[width=0.32\linewidth]{lamina_4}\\[2pt]
\includegraphics[width=0.32\linewidth]{lamina_2}\hfill
\includegraphics[width=0.32\linewidth]{lamina_3}\hfill
\includegraphics[width=0.32\linewidth]{lamina_8}
\caption{Algunas de las ilustraciones que acompañan a los combates.}
\end{figure}

\subsection{El turno de Morcar}

Cuando todos los héroes han jugado, el mapa ocupa la pantalla con el aviso
\pant{Es el turno de Morcar}. El ordenador mueve entonces, uno a uno, a los
monstruos que los héroes han descubierto, y los que alcanzan a un héroe lo
atacan. Cada ataque se resuelve en la pantalla de combate, igual que los de
los héroes. Al terminar, el juego lo anuncia y devuelve el turno al primer
héroe.

\pantalla[0.78]{turno_morcar}{El turno de Morcar se sigue sobre el mapa.}

\subsection{El mapa}

El icono del pergamino muestra el mapa de la mazmorra: sólo aparecen las
salas y pasillos que los héroes han explorado, con sus puertas, muebles,
trampas descubiertas y escaleras, además de la posición de cada héroe y de
los monstruos a la vista. Consultarlo no gasta el turno. Pulsa la acción
para volver a la sala.

El mismo mapa sirve para elegir el blanco de un ataque o de un hechizo.

\subsection{La magia}

Sólo el hechicero y el duende saben usar la magia. Cada hechizo puede
lanzarse \textbf{una sola vez} por aventura. Lanzar un hechizo sustituye al
ataque del turno y, como éste, puede hacerse antes o después de moverse.

Para lanzar un hechizo:

\begin{enumerate}
\item Abre el inventario (icono de la bolsa).
\item Pulsa sobre el hechizo en la lista \pant{Hechizos} y después sobre
  \pant{Salida}.
\item En el mapa, pulsa sobre el personaje o monstruo al que va dirigido.
  Un mensaje indica el resultado.
\end{enumerate}

\pantalla[0.7]{hoja_hechicero}{El inventario del hechicero, con sus nueve
hechizos.}

El blanco debe estar a la vista: lo están siempre quienes comparten sala
con el lanzador; en los pasillos, o entre salas distintas, hace falta una
línea recta sin muros, puertas cerradas ni otras figuras por medio. Si la
línea está cortada, el hechizo se desvía y se pierde. El lanzador siempre
puede hechizarse a sí mismo.

\subsubsection*{Hechizos de aire}
\begin{description}
\item[\pant{Genio}] Invoca a un genio que ataca una sola vez a cualquier
  enemigo con cinco dados de combate, o que abre una puerta.
\item[\pant{Tempestad}] Un torbellino envuelve a un monstruo o a un héroe,
  que pierde su próximo turno.
\item[\pant{Huracán}] El héroe elegido tirará el doble de dados en su
  próximo movimiento.
\end{description}

\subsubsection*{Hechizos de fuego}
\begin{description}
\item[\pant{Valor}] El héroe elegido tira dos dados más cada vez que ataca,
  hasta que deje de tener monstruos a la vista.
\item[\pant{Cólera de fuego}] Busca a un enemigo en cualquier lugar del
  tablero y le causa un punto de daño, salvo que saque un escudo con un
  dado.
\item[\pant{Bola de fuego}] Causa dos puntos de daño a un monstruo o héroe.
  La víctima tira dos dados y resta un punto por cada escudo.
\end{description}

\subsubsection*{Hechizos de tierra}
\begin{description}
\item[\pant{Piel de roca}] El héroe elegido tira dos dados más en defensa,
  hasta que resulte herido.
\item[\pant{Sanar heridas}] Devuelve hasta cuatro puntos de cuerpo
  perdidos.
\item[\pant{Atravesar muros}] En su próximo movimiento, el héroe elegido
  puede cruzar los muros, tantos como su movimiento le permita.
\end{description}

\subsubsection*{Hechizos de agua}
\begin{description}
\item[\pant{Dormir}] Sume a un monstruo o héroe en un sueño profundo, salvo
  que resista con sus puntos de mente. Dormido no puede defenderse. Puede
  despertar al comienzo de cada turno, y despierta si lo atacan.
\item[\pant{Velo de niebla}] En su próximo movimiento, el héroe elegido
  pasa sin ser visto por casillas ocupadas, y los monstruos no lo atacan.
\item[\pant{Aguas curativas}] Devuelve hasta cuatro puntos de cuerpo
  perdidos.
\end{description}

En el inventario algunos hechizos aparecen con el nombre abreviado
(\pant{Velo niebla}, \pant{Cruzar muros}, \pant{Sanar cuerpo}\ldots).

\subsection{El inventario, las pociones y los objetos}

El icono de la bolsa abre el inventario del héroe, dividido en
\pant{Pociones}, \pant{Armas}, \pant{Armadura}, \pant{Hechizos} (sólo el
duende y el hechicero) y \pant{Búsqueda tesoro}, que reúne los objetos
especiales hallados en las aventuras. Pulsa sobre un arma para empuñarla,
sobre una pieza de armadura para vestirla o sobre una poción para beberla;
lo que está en uso se muestra resaltado. \pant{Salida} cierra el
inventario.

\pantalla[0.7]{hoja_barbaro}{El inventario del bárbaro al empezar: todavía
vacío.}

\subsubsection*{Pociones}

Las pociones se encuentran buscando tesoros. Cada una sirve una sola vez.

\begin{table}[H]\centering\small
\begin{tabularx}{\linewidth}{@{}lL@{}}
\toprule
\textbf{Poción} & \textbf{Efecto}\\
\midrule
\pant{Curación}    & Devuelve al instante hasta cuatro puntos de cuerpo,
                     sin superar el máximo del héroe.\\
\pant{Fuerza}      & Dos dados más en el próximo ataque.\\
\pant{Resistencia} & Dos dados más en la próxima defensa.\\
\pant{Velocidad}   & Doble de dados en el movimiento del próximo turno.\\
\pant{Heroico}     & Permite atacar dos veces en la misma jugada.\\
\pant{Agua santa}  & Agua sagrada contra los muertos vivientes
                     (esqueletos, zombis y momias) de la sala.\\
\bottomrule
\end{tabularx}
\caption{Pociones.}
\end{table}

\subsubsection*{Tesoros de aventura}

Algunas aventuras esconden objetos únicos, que su dueño conserva de una
aventura a otra:

\begin{table}[H]\centering\small
\begin{tabularx}{\linewidth}{@{}lL@{}}
\toprule
\textbf{Objeto} & \textbf{Efecto}\\
\midrule
Armadura de Borin   & Cuatro dados de defensa, sin la lentitud de la
                      armadura de placas.\\
Vara del Recuerdo   & Permite lanzar dos hechizos en un mismo turno.\\
Talismán del Saber  & Aumenta en dos los puntos de mente de quien lo
                      lleva.\\
Azote de Orcos      & Espada de dos dados de ataque que golpea dos veces
                      cuando el enemigo es un orco.\\
Espada Espiritual   & Tres dados de ataque, cuatro contra los muertos
                      vivientes. Es la única arma capaz de herir al Señor
                      Brujo.\\
\bottomrule
\end{tabularx}
\caption{Tesoros de aventura.}
\end{table}

\subsection{La muerte}

Un héroe cuyos puntos de cuerpo llegan a cero muere y es retirado del
juego. Todo lo que llevaba, oro y armas, se pierde; los objetos únicos
vuelven al lugar donde se encontraron, de modo que, para recuperarlos, los
supervivientes tendrán que repetir la aventura correspondiente.

En la pantalla \pant{Crear personaje} el héroe caído figura como
\pant{Muerto}. Allí puede crearse un personaje nuevo que ocupe su lugar,
aunque tendrá que empezar desde cero.

\subsection{Entre aventuras}

Una aventura se supera al cumplir el objetivo que Mentor señaló en su
presentación y salir con vida de la mazmorra. Si no se consigue, puede
intentarse de nuevo: Morcar empieza siempre con todos sus monstruos.

Los héroes que sobreviven conservan su oro, su equipo y los tesoros de
aventura. También conservan sus heridas: los puntos de cuerpo perdidos no
vuelven solos, hay que recuperarlos con pociones o hechizos de curación. Es
el momento de pasar por la tienda antes de elegir la aventura siguiente, y
de apuntar el código de cada héroe (página~\pageref{sec:codigos}) si vas a
apagar el ordenador.

% ================================================================
\chapter{Los héroes}


\ficha{heroe_barbaro}{El bárbaro}{%
\begin{tabular}{@{}lc@{\qquad}lc@{}}
Puntos de cuerpo & 8 & Dados de ataque & 3\\
Puntos de mente & 2 & Dados de defensa & 2\\
\end{tabular}}{%
Nombre de fábrica: \textbf{Sigmar}. Es el guerrero más temible del grupo:
el que más golpes aguanta y el que con más fuerza los devuelve. No sabe
nada de magia, y su escasa fuerza mental lo hace vulnerable a los
hechizos, pero puede empuñar cualquier arma y vestir cualquier armadura.
Su lugar está en cabeza, abriendo las puertas y recibiendo a lo que
aguarde detrás.}

\ficha{heroe_enano}{El enano}{%
\begin{tabular}{@{}lc@{\qquad}lc@{}}
Puntos de cuerpo & 7 & Dados de ataque & 2\\
Puntos de mente & 3 & Dados de defensa & 2\\
\end{tabular}}{%
Nombre de fábrica: \textbf{Grugni}. Buen luchador, casi tan resistente
como el bárbaro, es además un experto en trampas: lleva consigo desde el
principio la caja de herramientas, con la que desactiva las trampas
descubiertas. Tampoco usa la magia. Con una buena arma comprada con el
primer oro se convierte en un combatiente formidable.}

\ficha{heroe_duende}{El duende}{%
\begin{tabular}{@{}lc@{\qquad}lc@{}}
Puntos de cuerpo & 6 & Dados de ataque & 2\\
Puntos de mente & 4 & Dados de defensa & 2\\
\end{tabular}}{%
Nombre de fábrica: \textbf{Ladril}. El elfo del juego de tablero, al que
esta versión llama \pant{Duende}, es guerrero y mago a la vez. Combate
como el enano, puede usar todas las armas y armaduras, y dispone además de
un grupo de tres hechizos, que elige al comenzar cada aventura.}

\ficha{heroe_hechicero}{El hechicero}{%
\begin{tabular}{@{}lc@{\qquad}lc@{}}
Puntos de cuerpo & 4 & Dados de ataque & 1\\
Puntos de mente & 6 & Dados de defensa & 2\\
\end{tabular}}{%
Nombre de fábrica: \textbf{Zoltan}. Débil en el cuerpo a cuerpo y con muy
pocos puntos de cuerpo, compensa su fragilidad con tres grupos de hechizos
---nueve en total--- y con la mente más fuerte del grupo. De la tienda
sólo le sirven el bastón, que le permite atacar en diagonal desde detrás
de un compañero, y la caja de herramientas. Hay que protegerlo: sus
hechizos, usados en el momento justo, deciden las aventuras.}

\begin{figure}[H]\centering
\includegraphics[width=0.48\linewidth]{heroe_barbaro_anim}\hfill
\includegraphics[width=0.48\linewidth]{heroe_enano_anim}\\[3pt]
\includegraphics[width=0.48\linewidth]{heroe_duende_anim}\hfill
\includegraphics[width=0.48\linewidth]{heroe_hechicero_anim}
\caption{Los cuatro héroes en movimiento.}
\end{figure}

% ================================================================
\chapter{Los enemigos}


Morcar dispone de ocho clases de criaturas. La tabla resume sus
características; el movimiento es el número de casillas que pueden
recorrer en un turno.

\begin{table}[H]\centering\small
\begin{tabular}{@{}llcccc@{}}
\toprule
\textbf{Enemigo} & \textbf{En pantalla} & \textbf{Mov.} & \textbf{Ataque}
 & \textbf{Defensa} & \textbf{Cuerpo}\\
\midrule
Goblin            & \pant{Duende}    & 10 & 2 & 1 & 1\\
Orco              & \pant{Orc}       & 8  & 3 & 2 & 1\\
Fimir             & \pant{Fimir}     & 6  & 3 & 3 & 1\\
Guerrero del Caos & \pant{Guerrero}  & 6  & 3 & 4 & 1\\
Esqueleto         & \pant{Esqueleto} & 6  & 2 & 2 & 1\\
Zombi             & \pant{Zombie}    & 4  & 2 & 3 & 1\\
Momia             & \pant{Momia}     & 4  & 3 & 4 & 1\\
Gárgola           & \pant{Hechicero} & 6  & 4 & 4 & 1\\
\bottomrule
\end{tabular}
\caption{Características de los enemigos (dados de ataque y de defensa).}
\end{table}

Recuerda que los monstruos sólo se defienden con una de las seis caras del
dado: incluso los mejor protegidos acaban cayendo.

\ficha{monstruo_goblin}{Goblin}{Movimiento 10 \quad Ataque 2 \quad Defensa 1}{%
Pequeños, cobardes y muy rápidos. El juego los llama \pant{Duende}, igual
que al héroe elfo: no los confundas. Uno solo es poca cosa, pero llegan
los primeros y suelen ir en grupo.}

\ficha{monstruo_orco}{Orco}{Movimiento 8 \quad Ataque 3 \quad Defensa 2}{%
La infantería de Morcar: fuertes, numerosos y siempre dispuestos a la
pelea. Pegan tan duro como el bárbaro. Los héroes se las verán con varios
de sus caudillos.}

\ficha{monstruo_fimir}{Fimir}{Movimiento 6 \quad Ataque 3 \quad Defensa 3}{%
Enormes reptiles de un solo ojo, armados con hachas y mazas. Más lentos
que los orcos, pero bastante más difíciles de abatir.}

\ficha{monstruo_guerrero_caos}{Guerrero del Caos}{Movimiento 6 \quad Ataque 3
\quad Defensa 4}{%
Campeones de Morcar enfundados en pesadas armaduras negras. Su defensa es
la mejor entre los vivos; suelen mandar a las demás criaturas y custodiar
los lugares importantes.}

\ficha{monstruo_esqueleto}{Esqueleto}{Movimiento 6 \quad Ataque 2 \quad Defensa 2}{%
Los muertos vivientes más comunes, alzados de sus tumbas por la magia del
Caos. Frágiles, pero incansables. Como todos los muertos vivientes, temen
al agua santa y a la Espada Espiritual.}

\ficha{monstruo_zombi}{Zombi}{Movimiento 4 \quad Ataque 2 \quad Defensa 3}{%
Cadáveres animados, lentos y torpes. Es fácil dejarlos atrás, aunque
resisten más de lo que su aspecto sugiere.}

\ficha{monstruo_momia}{Momia}{Movimiento 4 \quad Ataque 3 \quad Defensa 4}{%
Guardianas de las tumbas antiguas. Tan lentas como los zombis, pero
golpean como un orco y resisten como un guerrero del Caos.}

\ficha{monstruo_gargola}{Gárgola}{Movimiento 6 \quad Ataque 4 \quad Defensa 4}{%
La criatura más poderosa de Morcar: un demonio de piedra con el mejor
ataque y la mejor defensa del juego. La misma figura representa a los
hechiceros del Caos, y por eso el juego la llama \pant{Hechicero}.}

\subsection{Enemigos con nombre propio}

Algunas aventuras tienen un adversario principal, más peligroso que sus
congéneres y por cuya cabeza suele pagarse una recompensa:

\begin{description}
\item[Ulag] Señor de la guerra orco, responsable del secuestro de Sir
  Ragnar.
\item[Grak] Hijo de Ulag, que ha jurado vengar a su padre.
\item[Gulthor] Guerrero del Caos al mando de la banda que robó el oro del
  príncipe Magnus.
\item[Balur] El Mago del Fuego. Ningún hechizo de fuego puede dañarlo.
\item[El Señor Brujo] El \pant{Gran Brujo}, Rey de los Muertos. Es inmune a
  los hechizos, y de todas las armas sólo la Espada Espiritual puede
  herirlo.
\end{description}

En \textit{El retorno del Señor Brujo} aparecen además Bellthor, la gárgola
guardiana; Skulmar, capitán de la Legión Olvidada; Kessandria, la Reina
Bruja; los Jinetes Espectrales y los guerreros de la Guardia de la
Perdición, todos ellos más fuertes que las criaturas corrientes.

\subsection{El mobiliario de la mazmorra}

Las salas no están vacías: mesas, potros de tortura, bancos de alquimista,
tumbas, chimeneas, estanterías, armeros, cofres y tronos ocupan casillas y
cierran el paso. No se puede caminar sobre ellos, pero suelen señalar las
salas en las que merece la pena buscar.

\begin{figure}[H]\centering
\includegraphics[width=\linewidth]{mobiliario_a}\\[4pt]
\includegraphics[width=\linewidth]{mobiliario_b}\\[6pt]
\includegraphics[width=0.12\linewidth]{monstruo_cofre}\qquad
\includegraphics[width=0.12\linewidth]{monstruo_trono}
\caption{El mobiliario, el cofre y el trono.}
\end{figure}

% ================================================================
\chapter{Las aventuras de Morcar}


Éstas son las catorce aventuras del juego original, en el orden en que
conviene jugarlas. En el cartucho no hay que cargar nada: se eligen
directamente en la primera página de la lista.

\dospantallas{sala_q1}{sala_q6}{El comienzo de dos aventuras.}

\aventura{1}{El laberinto}
<<Ésta será vuestra prueba final antes de partir hacia las verdaderas
aventuras.>>

Mentor ha encerrado a sus discípulos en un laberinto y ha
soltado en él a varios monstruos.
\objetivo{Encontrar la salida. El primero que la alcance recibirá 100
monedas de oro para comprar equipo.}
\consejos{Los héroes empiezan separados, cada uno en una esquina. Es la
aventura ideal para aprender a moverse, abrir puertas y combatir. No hace
falta matar a todos los monstruos, y el hechicero hará bien en evitarlos.
El mapa te dirá por dónde has pasado ya.}

\aventura{2}{El rescate de Sir Ragnar}
Sir Ragnar, uno de los caballeros más poderosos del Emperador, ha sido
secuestrado. Lo retiene Ulag, señor de la guerra de los orcos.
\objetivo{Encontrar a Sir Ragnar y devolverlo sano y salvo a la escalera.
El príncipe Magnus pagará 200 monedas de oro a quien lo rescate. La
recompensa puede repartirse, pero no se pagará nada si Sir Ragnar muere
durante la huida.}
\consejos{Cuando los héroes dan con el prisionero, la alarma se extiende
por la fortaleza: todo el mapa queda al descubierto y los monstruos saben
dónde estáis. Sir Ragnar se une al grupo como un personaje más, pero está
débil, se mueve despacio y no puede luchar: escoltadlo. Buscad tesoros en
las salas ya limpias; revisad antes las trampas si veis un cofre.}

\aventura{3}{La guarida de los orcs}
El príncipe Magnus ha ordenado dar caza a Ulag, el responsable del
secuestro de Sir Ragnar.
\objetivo{Matar a Ulag. Quien lo consiga recibirá 100 monedas de oro, y
todo tesoro hallado en su fortaleza puede conservarse.}
\consejos{Ulag no está solo y pega más fuerte que un orco corriente:
atacadlo entre varios y reservad para él los hechizos de daño. Hay varios
tesoros escondidos ---entre ellos un arma--- y al menos un cofre con
trampa.}

\aventura{4}{El oro del príncipe Magnus}
Tres cofres han sido robados cuando viajaban hacia el Emperador. Los
ladrones son una banda de orcos escondida en las Montañas Negras, a las
órdenes de Gulthor, un guerrero del Caos.
\objetivo{Recuperar los cofres con \emph{todo} el oro y sacarlos de la
mazmorra. La recompensa es de 200 monedas de oro.}
\consejos{Un héroe que carga con un cofre queda marcado con \pant{Oro
Magnus} en su inventario y sólo tira un dado de movimiento. Despejad
primero el camino de vuelta, y que los héroes libres de carga protejan a
los porteadores.}

\aventura{5}{El laberinto Melar}
Hace mucho tiempo, el poderoso hechicero Melar creó un talismán que
acrecienta la comprensión de la magia de quien lo lleva. Temiendo que
cayera en manos de los aliados de Morcar, lo guardó en su laboratorio, en
el corazón de un laberinto protegido por trampas y guardianes mágicos. Se
dice que lo rondan los espíritus de quienes murieron buscándolo.
\objetivo{Encontrar el Talismán del Saber y salir con él.}
\consejos{Es una aventura de trampas y puertas secretas: buscad en cada
tramo de pasillo antes de avanzar y dejad que el enano vaya delante. No
todo se abre a la primera; antes del talismán hay que dar con otra cosa.
El talismán es perfecto para el hechicero o el duende.}

\aventura{6}{El legado orc}
Grak, el vástago de Ulag, ha jurado vengarse de quienes mataron a su
padre. Tras meses de búsqueda os ha tendido una emboscada y os ha
encerrado en sus mazmorras mientras decide vuestro castigo. Pero el
guardián duerme, y con un hueso de rata habéis forzado la cerradura de la
celda.
\objetivo{Recuperar vuestro equipo y escapar.}
\consejos{Empezáis sin armas ni armaduras, de modo que cualquier combate
es peligroso. Lo primero es encontrar la sala donde Grak guarda lo que os
quitó; hasta entonces, moveos juntos y evitad a los monstruos.}

\aventura{7}{El cazador de piedra}
Karlen, el mago personal del Emperador, ha desaparecido. El Emperador teme
que lo hayan asesinado, o que haya sucumbido a la seducción de la magia
del Caos.
\objetivo{Averiguar qué ha sido de Karlen y, si vive, ponerlo a salvo. A
la vuelta se pagarán 100 monedas de oro a cada héroe.}
\consejos{El título es una pista sobre lo que acecha en esta mazmorra. En
ella se esconden la Armadura de Borin, una de las mejores piezas del
juego, y el cofre de Karlen. Registrad las salas con muebles.}

\aventura{8}{Mago del fuego}
Los orcos de las Montañas Negras emplean magia de fuego en sus
incursiones. Se cree que los ayuda Balur, el Mago del Fuego, cuyos
hechizos no saben contrarrestar los magos del Emperador.
\objetivo{Penetrar en la guarida de Balur, bajo el Peñasco del Fuego Negro,
y destruirlo. El Emperador recompensará a cada héroe.}
\consejos{Ningún hechizo de fuego puede dañar a Balur, así que al repartir
la magia conviene que el hechicero cuente con otros grupos. En su guarida
se esconde la Vara del Recuerdo.}

\aventura{9}{Carrera contra el tiempo}
Un guía os ha conducido a una mazmorra de la que se dice que guarda un
gran secreto. Tras muchos corredores oscuros llegáis a una sala con tres
puertas; entonces el guía apaga su antorcha y, entre risas, se despide:
era una trampa.
\objetivo{Escapar con vida de ese agujero olvidado.}
\consejos{Empezáis sin saber dónde está la salida. No os separéis ni os entretengáis: los monstruos acuden desde
todas partes. Hay oro para quien se detenga a buscarlo, pero cada turno
cuenta.}

\aventura{10}{El castillo misterioso}
Hace mucho, el mago loco Ollar descubrió la entrada de una mina de oro y,
para protegerla, levantó sobre ella un castillo mágico, lleno de portales
encantados y guardado por monstruos atrapados en el tiempo.
\objetivo{Encontrar la entrada de la mina. Otros lo intentaron, y el
castillo los burló a todos.}
\consejos{Las puertas de este castillo no llevan adonde parece: al
cruzarlas, los héroes aparecen en otra sala, al azar. Es fácil que el
grupo se disperse. Quien encuentre el oro tendrá que decidir si carga con
él: con ese peso no se puede atacar ni defenderse, aunque siempre es
posible soltarlo desde el inventario.}

\aventura{11}{El baluarte de Caos}
Las tierras del este sufren el saqueo de orcos y goblins, que se refugian
en una fortaleza subterránea, el Baluarte del Caos, al mando de un pequeño
grupo de guerreros del Caos.
\objetivo{Abrirse paso y acabar con todos los monstruos. Se pagan 10
monedas por cada goblin, 20 por cada orco y 30 por cada fimir o guerrero
del Caos.}
\consejos{Es la aventura con más combates: entrad con el mejor equipo que
podáis pagar y con todos los héroes. Avanzad sala a sala, sin abrir una
puerta nueva hasta haber limpiado la anterior. Hay una gárgola que parece
una estatua; mientras nadie toque lo que guarda, seguirá siéndolo.}

\aventura{12}{Barak Tor}
Se avecina la guerra contra los orcos del este, y el Emperador necesita
unir a los reinos menores. Para ello debe encontrar la Estrella del Oeste,
la joya que llevaron los reyes de la leyenda y el propio Rogar. La gema
yace en Barak Tor, el túmulo donde reposa el Señor Brujo.
\objetivo{Encontrar la Estrella del Oeste y salir con ella. Quien la
entregue recibirá 200 monedas de oro.}
\consejos{Barak Tor es un lugar de muertos vivientes: esqueletos, zombis
y momias, lentos pero difíciles de abatir. Y cuidado con lo que
despertáis: hay sellos que es mejor no romper, y derrumbes que pueden
cortar la retirada.}

\aventura{13}{La espada espiritual}
Habéis despertado al Señor Brujo, el Rey de los Muertos, un poderoso
siervo de Morcar. Hay que destruirlo antes de que lance su ejército de
muertos vivientes contra las fuerzas del Emperador, y sólo un arma puede
herirlo: la Espada Espiritual, forjada por los enanos de las Montañas del
Fin del Mundo y templada en la fuente élfica de Lebin. La espada yace en
un antiguo templo en ruinas.
\objetivo{Encontrar la Espada Espiritual.}
\consejos{El templo está bien guardado. La espada no es el único tesoro
que esconde. Entregádsela a un buen luchador: contra los muertos vivientes
tira cuatro dados, y sin ella la aventura siguiente no puede ganarse.}

\aventura{14}{Regreso a Barak Tor}
Con la Espada Espiritual en vuestro poder debéis volver a Barak Tor. El
Emperador ha partido al encuentro de los orcos del este en el Paso del
Fuego Negro; si fracasáis, el Señor Brujo atacará a sus tropas por la
retaguardia con su ejército de muertos vivientes y nada impedirá que las
fuerzas del Caos invadan la tierra.
\objetivo{Derrotar al Señor Brujo.}
\consejos{El Señor Brujo es inmune a los hechizos y a las armas
corrientes: sólo el portador de la Espada Espiritual puede herirlo. Los
demás deben protegerlo y mantener a raya a la escolta de muertos
vivientes. Guardad los hechizos de curación para él.}

% ================================================================
\chapter{El retorno del Señor Brujo}


La segunda página de la lista de aventuras contiene las diez de la
expansión, que continúan la historia donde la dejó \textit{Regreso a Barak
Tor}. Son más largas y difíciles que las anteriores, y es todavía más
importante jugarlas en orden y con héroes bien equipados.

\pantalla[0.7]{pergamino_exp}{La presentación de la primera aventura de la
expansión.}

\subsection*{Un mensaje de Mentor}

<<Dignos compañeros, hay que actuar deprisa: ¡el Señor Brujo no ha
muerto!>> Logró huir a Kalos, la ciudad caída. Aquella hermosa ciudad fue
destruida hace muchos años por Morcar, que la entregó al Señor Brujo; y
bajo sus ruinas levantó éste una gran fortaleza, pues sus criaturas no
soportan la luz del día. En cuanto recobre las fuerzas reunirá a su
ejército de muertos vivientes y atacará al Imperio. Sus guerreros no
conocen el miedo y luchan hasta ser destruidos.

Muchos días de marcha por las abrasadas Llanuras de la Muerte os han
traído hasta las ruinas de Kalos. Durante todo el viaje habéis sentido una
presencia inquietante: desde su alto trono, el Señor Brujo ve a cuantos se
acercan. Si no se le detiene pronto se le unirán Skulmar, capitán de la
Legión Olvidada, y Kessandria, la Reina Bruja; y se dice que, si ambos se
sientan a su lado, los huesos de todos los guerreros caídos se alzarán
bajo el Estandarte Negro.

\subsection*{Novedades de la expansión}

\begin{itemize}
\item Los héroes no siempre empiezan en la escalera: a menudo entran por
  una puerta en el borde de la mazmorra, y la salida puede ser otra puerta
  distinta, que hay que encontrar.
\item Entre una aventura y otra los héroes recuperan todos sus puntos de
  cuerpo y de mente.
\item Hay enemigos con nombre propio mucho más fuertes que los corrientes,
  nuevas trampas y salas con comportamientos mágicos.
\end{itemize}

\dospantallas{sala_q10}{sala_q13}{Los héroes, preparados para descender una
vez más.}

\aventura{1}{La puerta de la perdición}
La entrada a la fortaleza no es difícil de encontrar: nadie se ha
molestado en ocultarla. La puerta exterior está abierta y hay señales de
que muchas criaturas la han cruzado en los últimos días.
\objetivo{Abrirse camino por los túneles hasta la Puerta de la
Perdición.}
\consejos{El enemigo ha tenido tiempo de preparar la bienvenida: hay
puertas falsas que no se abren y salas cuyos ocupantes aguardan
emboscados y salen todos a la vez. Merece la pena registrar las salas:
hay oro, y algo más, para quien busque.}

\aventura{2}{Las salas frías}
Estos pasadizos helados son el lugar de reposo de los Jinetes
Espectrales, que antaño guiaron a la batalla a la caballería de
esqueletos del Señor Brujo. Ahora yacen soñando con glorias pasadas, a la
espera de que el Estandarte Negro vuelva a llamarlos.
\objetivo{Encontrar la escalera de caracol para seguir adelante.}
\consejos{Los Jinetes Espectrales parecen esqueletos, pero atacan y se
defienden mucho mejor. Hay una sala giratoria de
la que no se sabe por qué puerta se saldrá, y un cofre protegido por una
aguja envenenada: revisad las trampas antes de abrirlo. No todas las
tumbas deben profanarse.}

\aventura{3}{Los pasadizos silenciosos}
Estos pasadizos llevan en silencio más de mil años. Las últimas voces que
resonaron en ellos fueron las del rey Agrain y su pueblo, cuando corrían a
defender la ciudad. Se dice que aquí actúa una magia antigua y terrible.
\objetivo{Encontrar la puerta de salida.}
\consejos{Una niebla mortal, aliento del Caos, recorre los corredores y
daña a todo el que no sea maligno. No os quedéis en su camino. En las
salas hay un cofre con trampa y algunas armas arrojadizas mágicas.}

\aventura{4}{Las salas de la visión}
El rey Agrain fue uno de los grandes reyes anteriores al Caos, guerrero
temido y legislador justo. También jugó con la magia, y así creó las Salas
de la Visión, un laberinto de salas mágicas cuyo secreto está en las
Llaves de Agrain.
\objetivo{Encontrar las Llaves de Agrain y alcanzar la puerta de salida.}
\consejos{Cuatro salas mágicas, guardada cada una por una momia, protegen
el laboratorio secreto del rey. Al caer una momia, las puertas cambian.
Las llaves están en el laboratorio, y son valiosas por más de un motivo.}

\aventura{5}{La puerta de Bellthor}
Aquí empieza el verdadero dominio del Señor Brujo: la fortaleza
subterránea desde la que planea la destrucción del Imperio. Loretome habla
de un guardián terrible, Bellthor, que vigila la puerta.
\objetivo{Enfrentarse a Bellthor. Sólo juntos podréis vencerlo.}
\consejos{Bellthor es una gárgola con un aliento venenoso que ataca la
mente, y no se mueve hasta tener a todos los héroes delante. Entrad
juntos, con todos los puntos de cuerpo posibles, y recoged antes el oro de
los cofres. La victoria tendrá un precio inesperado.}

\aventura{6}{Las salas de los muertos}
El hechicero y el enano despiertan en una celda, todavía aturdidos por el
gas venenoso. Sois prisioneros del Señor Brujo, que reúne aquí su ejército
inmundo; el duende y el bárbaro han sido llevados a un lugar secreto.
\objetivo{Rescatar a los compañeros y escapar por la escalera de caracol.}
\consejos{Al principio sólo juegan el hechicero y el enano. Conservan sus
armas, objetos y hechizos, pero el oro ha desaparecido. El bárbaro y el
duende están encadenados y sin equipo: tras liberarlos hay que encontrar
el cofre donde se guarda.}

\aventura{7}{La legión olvidada}
Habéis escapado del Señor Brujo, pero no tardará en seguiros. Las salas
que ahora cruzáis guardan a la Legión Olvidada, la élite de su ejército,
al mando de Skulmar: guerreros con las cicatrices de mil batallas, cada
una de ellas una victoria del Caos.
\objetivo{Encontrar la puerta de salida.}
\consejos{La Legión aguarda inmóvil, encerrada en un hielo mágico que
ningún arma rompe\ldots{} hasta que se abre la puerta equivocada. Skulmar
es más fuerte que cualquier guerrero del Caos y, si la lucha se tuerce,
intentará huir; si cae, deja tras de sí un botín interesante.}

\aventura{8}{La ciudad prohibida}
Ningún hombre ha entrado jamás en la Ciudad Prohibida, y ni siquiera
Loretome sabe decir gran cosa de lo que aguarda. Se dice que sus pasadizos
están llenos de muertos vivientes, los siervos sin mente de Morcar,
condenados a vagar por ella para siempre.
\objetivo{Atravesar la ciudad y encontrar la salida.}
\consejos{Los muertos vivientes no se acaban: aparecen otros nuevos turno
tras turno, de modo que no tiene sentido quedarse a pelear. La ciudad,
además, se desmorona. Hay tesoros para quien se arriesgue, alguno de ellos
envenenado.}

\aventura{9}{La última puerta}
Ésta es la última puerta antes de la corte del Señor Brujo. Los pasadizos
que hay tras ella están guardados por los guerreros del Caos de la Guardia
de la Perdición, más poderosos que cuantos habéis combatido, y en ellos
tiene sus aposentos Kessandria, la Reina Bruja.
\objetivo{Encontrar la escalera de caracol que lleva ante el Señor Brujo.}
\consejos{Los guardias de la Perdición atacan con cuatro dados y se
defienden con seis. Kessandria domina la magia
del Caos y sólo teme a los hechizos de fuego: guardadlos para ella.}

\aventura{10}{La corte del Señor Brujo}
El Señor Brujo ha observado vuestro avance. Ha intentado destruiros una y
otra vez, y otras tantas lo habéis frustrado. Ahora ha huido a su sala del
trono.
\objetivo{Librar allí la batalla final. Esta vez no debe haber escapatoria
para él.}
\consejos{El Señor Brujo es ahora más fuerte que en Barak Tor, conoce los
hechizos del Caos y lo escoltan guerreros de la Guardia de la Perdición.
Sólo pueden herirlo la Espada Espiritual, los hechizos Cólera de fuego y
Bola de fuego y las dagas arrojadizas mágicas. No ataquéis a las estatuas
de orcos que cierran los corredores: no sufren daño, y las armas se
rompen contra ellas.}

\subsection*{Conclusión}

Si el Señor Brujo cae, su ejército de muertos vivientes seguirá dormido
bajo la tierra y la magia que lo gobernaba quedará debilitada. Pero Mentor
advierte de que la victoria no debe ser excusa para el descanso: los
siervos que escaparon buscarán venganza, y Skulmar, a quien se creía
muerto, ha huido de las ruinas de Kalos con muchos de los libros de
hechizos de su amo. Por ahora, sin embargo, os habéis ganado un respiro:
el Emperador os espera en el Gran Palacio.

% ================================================================
